% =======================================================================================
% BRAC University — Beamer presentation template
%
% Built on Beamer and its Madrid theme, both by Till Tantau (Madrid, 2007), now maintained
% by Joseph Wright, Vedran Miletic and contributors. Beamer is licensed under the LPPL and
% the GNU GPL. What is added here is the BRAC University colouring, the institute block,
% the one-file-per-section layout and the bibliography wiring.
%
% A clean starting point for a thesis defence, a seminar or a conference talk.
%
% HOW TO USE
%   1. Change the title, author and supervisor below.
%   2. Replace img/bracu_logo.png if you are at another institution, and change the three
%      colours at the top of config.tex to match your own logo.
%   3. Write each section in its own file under files/ and \input it below. One file per
%      section keeps the deck manageable and makes version control sane.
%   4. Put references in ref.bib and cite them with \cite{key}.
%
% COMPILING
%   Overleaf : set this file as the main document, compiler pdfLaTeX, references BibTeX.
%   Locally  : latexmk -pdf main.tex
%
% WHAT IS WHERE
%   config.tex   packages, colours, the institute block, the per-section contents slide
%   files/       one .tex per section, in the order they appear below
%   img/         logo and figures
%   ref.bib      bibliography
% =======================================================================================

\documentclass[11pt]{beamer}

%%%%%%%% theme %%%%%%%%
% Madrid gives a title bar and a footer with author, title, date and slide number.
% Other themes worth trying: Frankfurt, Singapore, CambridgeUS, metropolis (needs install).
\mode<presentation> {
\usetheme{Madrid}
%\usecolortheme{albatross}
}

% ---------------------------------------------------------------------------------------
% Packages, colours and the title-block institute.
%
% Everything that is not slide content lives here, so the files in files/ stay readable.
% Change the colours in one place below and the whole deck follows.
% ---------------------------------------------------------------------------------------

\usepackage[english]{babel}
\usepackage[utf8]{inputenc}
\usepackage{graphicx}
\usepackage{booktabs}     % proper table rules: \toprule \midrule \bottomrule
\usepackage{amsmath}
\usepackage{amssymb}
\usepackage{amsfonts}
\usepackage{comment}      % \begin{comment} ... \end{comment} to park a draft slide

% --- Colours ---------------------------------------------------------------------------
% None. The deck uses the Madrid theme's own palette, which is well balanced and prints
% legibly, so nothing is overridden here.
%
% If you want a different look, the least invasive route is a colour theme — uncomment one
% in main.tex, for example \usecolortheme{albatross}, \usecolortheme{seahorse} or
% \usecolortheme{dolphin}.
%
% To match an institution's own colours instead, add something like the following and
% adjust the hex values. Pick them from the logo file with any image editor's colour
% picker rather than guessing. One warning: Madrid draws the title on a coloured panel, so
% a title foreground in the same colour renders an invisible title.
%
%   \usepackage{xcolor}
%   \definecolor{primary}{HTML}{1F58A9}
%   \setbeamercolor{structure}{fg=primary}
%   \setbeamercolor{frametitle}{bg=primary,fg=white}
%   \setbeamercolor{title}{bg=primary,fg=white}
%   \setbeamercolor{block title}{bg=primary,fg=white}

% The little navigation arrows are never used in a talk and add visual noise.
\setbeamertemplate{navigation symbols}{}
\setbeamertemplate{caption}[numbered]

% --- Title block --------------------------------------------------------------------------
% Appears on the title slide. Replace the logo in img/ and the two lines below.
\institute[BRAC University]
{
% Keep the logo this size or smaller. The title block holds logo, institute, title,
% subtitle, author and date; at scale 0.35 a two-line subtitle pushes it off the slide.
\vspace*{-0.5cm}
\begin{figure}[h]
 \centering
 \includegraphics[scale=0.28]{img/bracu_logo.png}
\end{figure}
BRAC University\\
Department of Computer Science \& Engineering
}
\date{\today}

% --- A contents slide before each section -------------------------------------------------
% Useful in a long talk, tiring in a short one: each section costs a slide. Comment this
% block out if your deck has many short sections.
\AtBeginSection[]
{
\begin{frame}
\frametitle{Contents}
\tableofcontents[currentsection]
\end{frame}
}


%%%%%%%% title %%%%%%%%
% The short title in [ ] appears in the footer of every slide — keep it brief.
\title[Short title]{Your Presentation Title}
\subtitle{An optional subtitle that says what the work is about}

%%%%%%%% author %%%%%%%%
% The short name in [ ] also goes in the footer.
\author[Your Name]{
Your Name\\\vspace{0.15cm}
\small Supervisor: Supervisor Name}

\begin{document}

% Title slide and the opening contents slide.

\begin{frame}
\titlepage
\end{frame}

\begin{frame}
\frametitle{Contents}
\tableofcontents
\end{frame}


%%%%%%%% slides %%%%%%%%
% Add, remove or reorder sections here. Each \section creates a contents slide
% (see config.tex) and an entry in the table of contents.
\section{Introduction}
% Section 1 — Introduction.
% This file demonstrates the building blocks you will use most: blocks, lists, emphasis
% and citations. Delete the examples and write your own.

\begin{frame}{Motivation}

% A block with a heading. Use it to group one idea per block.
\begin{block}{The problem}
\begin{itemize}
    \item State the problem in one line that a non-specialist understands
    \item Say who is affected, and how much
    \item Resist the urge to explain your method here
\end{itemize}
\end{block}

% A block with no heading, written as {}, is a quiet container for supporting points.
\begin{block}{}
\begin{itemize}
    \item Supporting detail, with a citation~\cite{W3021074122}
    \item A second point, \textbf{emphasised} where it matters
\end{itemize}
\end{block}

\end{frame}


% [allowframebreaks] lets a long frame continue onto a second slide automatically.
% Beamer appends "I", "II" to the title. Useful while drafting; tighten before presenting.
\begin{frame}[allowframebreaks]{Background}

\begin{block}{What is already known}
\begin{itemize}
    \item One finding per line, each attached to its source~\cite{W3135286407}
    \item Keep lines short enough to read from the back of the room
\end{itemize}
\end{block}

\begin{block}{What is not known}
\begin{itemize}
    \item The gap your work fills, stated as plainly as you can manage
\end{itemize}
\end{block}

\end{frame}


\section{Methodology}
% Section 2 — Methodology.
% Demonstrates: two columns, a figure, and an equation.

\begin{frame}{Approach}

% Two columns. Handy for text beside a figure.
\begin{columns}[T]
\begin{column}{0.52\textwidth}
\begin{block}{Steps}
\begin{itemize}
    \item Data collection
    \item Pre-processing
    \item Model fitting
    \item Validation
\end{itemize}
\end{block}
\end{column}

\begin{column}{0.44\textwidth}
\begin{block}{Why this way}
\begin{itemize}
    \item A reason, not a restatement
    \item The alternative you rejected
\end{itemize}
\end{block}
\end{column}
\end{columns}

\end{frame}


\begin{frame}{A figure}

% Put image files in img/. scale= or width= both work; width relative to \linewidth
% travels better between screens and projectors.
\begin{figure}
\centering
\includegraphics[width=0.45\linewidth]{img/bracu_logo.png}
\caption{Replace this with your own figure. Captions are numbered automatically.}
\end{figure}

\end{frame}


\begin{frame}{An equation}

\begin{block}{}
One displayed equation per slide is usually the limit an audience can absorb:
\begin{equation}
    \hat{y} = \sigma\!\left(\beta_0 + \sum_{i=1}^{n} \beta_i x_i\right)
\end{equation}
Define every symbol you use, in words, immediately after it.
\end{block}

\end{frame}


\section{Results}
% Section 3 — Results.
% Demonstrates a table using booktabs, which looks far better than the default rules.

\begin{frame}{Results}

\begin{table}
\centering
\begin{tabular}{lrr}
\toprule
\textbf{Model} & \textbf{Accuracy} & \textbf{AUC} \\
\midrule
Baseline       & 0.71 & 0.74 \\
Proposed       & 0.83 & 0.88 \\
\bottomrule
\end{tabular}
\caption{Keep tables small. If it needs a scrollbar, it needs a figure instead.}
\end{table}

\begin{block}{What the numbers mean}
\begin{itemize}
    \item Say what the reader should take away before they work it out themselves
    \item Note the comparison that makes the number meaningful
\end{itemize}
\end{block}

\end{frame}


\section{Conclusion}
% Section 4 — Conclusion.

\begin{frame}{Conclusion}

\begin{block}{Contributions}
\begin{itemize}
    \item What exists now that did not exist before your work
\end{itemize}
\end{block}

\begin{block}{Limitations}
\begin{itemize}
    \item Name them yourself — an examiner who finds one you missed trusts the rest less
\end{itemize}
\end{block}

\begin{block}{Next}
\begin{itemize}
    \item The question you would answer with another six months
\end{itemize}
\end{block}

\end{frame}


\begin{frame}
\begin{center}
\Large Thank you \\[0.5cm]
\normalsize Questions?
\end{center}
\end{frame}


%%%%%%%% references %%%%%%%%
% \nocite{*} would list everything in ref.bib, cited or not. Leave it off to list only
% what you actually cite.
\begin{frame}[allowframebreaks]{References}
\footnotesize
\bibliographystyle{unsrt}
\bibliography{ref}
\end{frame}

%%%%%%%% closing slide %%%%%%%%
% Repeating the title slide gives the audience something to look at during questions.
\begin{frame}
\titlepage
\end{frame}

\end{document}
